

\ifdefined\maindoc
% Under the thesis build, main.tex owns the preamble and the \chapter/\label lines.
\else
\documentclass[12pt]{report}
\usepackage[a4paper,margin=1in]{geometry}
\usepackage[utf8]{inputenc}
\usepackage[T1]{fontenc}
\usepackage{amsmath,amssymb,amsthm}
\usepackage{booktabs}
\usepackage{setspace}
\usepackage{graphicx}
\usepackage{longtable}
\usepackage{array}
\usepackage{makecell}
\usepackage{enumitem}
\usepackage{placeins}
\usepackage{hyperref}
\usepackage[toc,page]{appendix}
\onehalfspacing
\begin{document}
\title{Appendix: Case-Study Network Specifications}
\fi

% ---------------------------------------------------------------------------
% CONTENT NOTE (delete once populated)
%
% This appendix holds the canonical, machine-checkable definition of every
% network used as a case study anywhere in the thesis, so each chapter can
% cite one authoritative source instead of re-describing a network in prose.
%
% Expected content, one section per network:
%   - a short identifier and which chapter(s) use it
%     (e.g. the flagship 43-node/107-edge flow benchmark, the 8-node nested
%      diamond network, the water accumulation schedule, the medical drone
%      network, the RESS multi-level reliability network, the grid 5x5 mesh)
%   - node list: id, human-readable label, role (source / sink / junction / ...)
%   - edge list: (u, v) pairs
%   - the analysis inputs actually used: edge/node capacities for the flow
%     runs, activity durations for the CPM runs, component and handover
%     probabilities (point / interval / p-box) for the reliability runs
%   - provenance: the .EDGES / .json file the numbers come from, and the
%     validation script that regenerates any derived figures
%
% Keep this to data and provenance. Interpretation stays in the chapters.
% ---------------------------------------------------------------------------

\emph{Appendix in preparation.} See the content note in the source file for the intended
structure. Cross-reference target: \texttt{app:network-specs}.
